% Generated by generate.py from src/ch-threads.tex: edit the source, not this file.
\chapter{Threads Between the Tomes}\label{ch:threads}

A theorem proved once is used many times. This chapter follows the references from one tome to another: every place
where a tome writes ``Tome~$N$, Theorem~$x.y$'' or ``Tome~$N$, Chapter~$k$''. Each reference was checked against the
sources when this book was generated: the result it names exists, and has the kind the reference says.

\section{Who cites whom}

Figure~\ref{fig:citations} in Chapter~\ref{ch:map} shows the counts as a table; here they are tome by tome, in words.

% Generated by generate.py: do not edit.
\paragraph{Tome 0} cites Tome~1 (2).
\paragraph{Tome 1} cites Tome~0 (12).
\paragraph{Tome 2} cites Tome~0 (12), Tome~1 (13).
\paragraph{Tome 3} cites Tome~0 (21), Tome~2 (1), Tome~4 (1), Tome~C1 (1).
\paragraph{Tome 4} cites Tome~0 (18), Tome~1 (8), Tome~2 (3), Tome~3 (7), Tome~5 (1).
\paragraph{Tome 5} cites Tome~1 (24), Tome~2 (6), Tome~3 (6), Tome~4 (68).
\paragraph{Tome 6} cites Tome~0 (3), Tome~1 (5), Tome~2 (1), Tome~3 (7), Tome~4 (3), Tome~7 (1).
\paragraph{Tome 7} cites Tome~1 (7), Tome~2 (5), Tome~3 (5), Tome~4 (41), Tome~5 (6), Tome~6 (66), Tome~S1 (1).
\paragraph{Tome 8} cites Tome~1 (1), Tome~2 (3), Tome~3 (1), Tome~5 (5), Tome~6 (15), Tome~7 (14), Tome~12 (2), Tome~S1 (4), Tome~S2 (6).
\paragraph{Tome 9} cites Tome~0 (1), Tome~1 (1), Tome~2 (47), Tome~3 (5), Tome~5 (1), Tome~6 (8), Tome~7 (3), Tome~8 (17), Tome~12 (4), Tome~S1 (4).
\paragraph{Tome 10} cites Tome~1 (2), Tome~3 (73), Tome~4 (6), Tome~5 (3), Tome~6 (32), Tome~7 (6), Tome~C1 (1).
\paragraph{Tome 11} cites Tome~3 (45), Tome~4 (3), Tome~5 (2), Tome~6 (18), Tome~7 (4), Tome~8 (7), Tome~10 (24), Tome~C1 (8), Tome~S1 (1), Tome~S3 (4).
\paragraph{Tome 12} cites Tome~0 (1), Tome~1 (35), Tome~2 (42), Tome~4 (34), Tome~5 (20), Tome~6 (8).
\paragraph{Tome 13} cites Tome~3 (1), Tome~4 (13), Tome~5 (62), Tome~7 (5), Tome~10 (3), Tome~11 (4), Tome~12 (9), Tome~C1 (2).
\paragraph{Tome 14} cites Tome~1 (6), Tome~2 (2), Tome~3 (9), Tome~4 (9), Tome~5 (52), Tome~6 (6), Tome~8 (1), Tome~10 (16), Tome~11 (8), Tome~12 (6), Tome~13 (31), Tome~15 (1).
\paragraph{Tome 15} cites Tome~0 (1), Tome~2 (2), Tome~3 (19), Tome~5 (3), Tome~9 (11), Tome~10 (3), Tome~11 (5), Tome~12 (2), Tome~13 (2), Tome~14 (14), Tome~C1 (3).
\paragraph{Tome C1} cites Tome~0 (4), Tome~1 (9), Tome~2 (13), Tome~3 (14), Tome~4 (17), Tome~6 (3), Tome~12 (1), Tome~C2 (1), Tome~S1 (7).
\paragraph{Tome C2} cites Tome~0 (1), Tome~1 (14), Tome~2 (8), Tome~3 (8), Tome~4 (7), Tome~C1 (27), Tome~S1 (5).
\paragraph{Tome C3} cites Tome~1 (5), Tome~2 (6), Tome~3 (2), Tome~4 (1), Tome~6 (2), Tome~C1 (31), Tome~C2 (31).
\paragraph{Tome C4} cites Tome~1 (1), Tome~3 (17), Tome~C2 (2), Tome~C3 (8).
\paragraph{Tome C5} cites Tome~3 (1), Tome~6 (1), Tome~10 (4), Tome~C1 (3), Tome~S1 (13).
\paragraph{Tome S1} cites Tome~0 (1), Tome~1 (2), Tome~2 (4), Tome~3 (1), Tome~4 (3), Tome~5 (1), Tome~6 (3), Tome~7 (2), Tome~12 (1), Tome~C1 (3), Tome~C5 (2), Tome~S2 (2).
\paragraph{Tome S2} cites Tome~0 (1), Tome~1 (1), Tome~2 (1), Tome~12 (4), Tome~S1 (1).
\paragraph{Tome S3} cites Tome~0 (1), Tome~1 (7), Tome~2 (11), Tome~3 (5), Tome~4 (14), Tome~5 (50), Tome~7 (2), Tome~12 (17), Tome~13 (5), Tome~14 (1).
\paragraph{Tome P1} cites Tome~0 (156), Tome~1 (179), Tome~2 (131), Tome~3 (187), Tome~4 (200), Tome~5 (174), Tome~6 (1), Tome~C1 (5).
\paragraph{Tome P2} cites Tome~0 (2), Tome~1 (14), Tome~3 (7), Tome~4 (28), Tome~5 (7), Tome~6 (62), Tome~7 (33), Tome~10 (21), Tome~13 (4), Tome~S1 (2).
\paragraph{Tome P3} cites Tome~0 (1), Tome~1 (2), Tome~2 (1), Tome~3 (19), Tome~4 (11), Tome~5 (5), Tome~6 (8), Tome~7 (4), Tome~9 (2), Tome~10 (32), Tome~12 (1).


\section{The results the series relies on most}

These are the numbered results cited by the largest number of other tomes, with the tome where each is proved and the
tomes that use it.

% Generated by generate.py: do not edit.
\begin{results}
\item[3] Theorem~4.11 (Borel--Cantelli lemmas), in \emph{Inequalities, Large Numbers and Entropy}; cited by Tome~5, Tome~10, Tome~14, Tome~C4, Tome~P1, Tome~P2.
\item[4] Theorem~10.3 (Euler product), in \emph{The Zeta Function}; cited by Tome~5, Tome~7, Tome~12, Tome~S3, Tome~P1, Tome~P2.
\item[5] Theorem~1.5 (Mertens, 1874), in \emph{Mertens's Theorems}; cited by Tome~10, Tome~13, Tome~14, Tome~S3, Tome~P1, Tome~P3.
\item[6] Theorem~11.7 (Prokhorov), in \emph{Measures on Metric Spaces}; cited by Tome~7, Tome~10, Tome~11, Tome~14, Tome~P2, Tome~P3.
\item[3] Theorem~6.15, in \emph{Convergence and the Central Limit Theorem}; cited by Tome~10, Tome~11, Tome~C5, Tome~P1, Tome~P2.
\item[3] Theorem~6.8 (uniqueness and inversion), in \emph{Convergence and the Central Limit Theorem}; cited by Tome~5, Tome~10, Tome~14, Tome~P1, Tome~P3.
\item[4] Proposition~10.6, in \emph{The Zeta Function}; cited by Tome~5, Tome~7, Tome~12, Tome~14, Tome~S3.
\item[4] Theorem~6.13 (Weyl's criterion, 1916), in \emph{Fourier Series}; cited by Tome~5, Tome~13, Tome~14, Tome~P1, Tome~P2.
\item[4] Theorem~8.9 (Poisson summation), in \emph{The Fourier Transform and Poisson Summation}; cited by Tome~7, Tome~12, Tome~13, Tome~P1, Tome~P3.
\item[4] Theorem~9.15 (limits and integrals), in \emph{Complex Analysis for Computing}; cited by Tome~5, Tome~7, Tome~12, Tome~13, Tome~P2.
\item[4] Theorem~9.18 (zeros and the identity theorem), in \emph{Complex Analysis for Computing}; cited by Tome~5, Tome~7, Tome~10, Tome~P2, Tome~P3.
\item[5] Theorem~5.2 (Hadamard, de la Vallée Poussin, 1896), in \emph{The Prime Number Theorem}; cited by Tome~7, Tome~13, Tome~S3, Tome~P1, Tome~P2.
\item[5] Theorem~5.5 (the prime number theorem), in \emph{The Prime Number Theorem}; cited by Tome~7, Tome~14, Tome~S3, Tome~P1, Tome~P2.
\item[1] Theorem~2.4 (Chebyshev-type bounds), in \emph{Primes}; cited by Tome~5, Tome~C2, Tome~C4, Tome~P1.
\item[1] Exercise~6.7, in \emph{Quadratic Residues and Reciprocity}; cited by Tome~4, Tome~12, Tome~P1, Tome~P3.
\item[1] Theorem~7.3, in \emph{Sums of Squares}; cited by Tome~7, Tome~14, Tome~C1, Tome~P1.
\item[2] Theorem~11.14 (Fourier analysis on a finite abelian group), in \emph{Representations and Characters}; cited by Tome~4, Tome~7, Tome~9, Tome~P1.
\item[2] Theorem~9.6 (Galois, Moore), in \emph{Finite Fields}; cited by Tome~12, Tome~C1, Tome~C2, Tome~S3.
\item[3] Theorem~4.6 (Chernoff bounds), in \emph{Inequalities, Large Numbers and Entropy}; cited by Tome~5, Tome~10, Tome~C2, Tome~C3.
\item[3] Theorem~4.9 (Hoeffding's inequality), in \emph{Inequalities, Large Numbers and Entropy}; cited by Tome~5, Tome~C2, Tome~C3, Tome~C4.
\item[3] Theorem~6.10 (Lindeberg, 1922), in \emph{Convergence and the Central Limit Theorem}; cited by Tome~10, Tome~15, Tome~P1, Tome~P3.
\item[4] Theorem~10.8 (Riemann, 1859), in \emph{The Zeta Function}; cited by Tome~5, Tome~12, Tome~13, Tome~S3.
\item[4] Theorem~8.6 (Fourier inversion), in \emph{The Fourier Transform and Poisson Summation}; cited by Tome~12, Tome~13, Tome~14, Tome~P1.
\item[4] Theorem~9.36, in \emph{Complex Analysis for Computing}; cited by Tome~5, Tome~10, Tome~12, Tome~14.
\item[5] Theorem~3.5, in \emph{Dirichlet Characters and \texorpdfstring{$L$}{L}-Functions}; cited by Tome~12, Tome~13, Tome~S3, Tome~P1.
\item[5] Theorem~6.1 (Stirling's formula), in \emph{The Zeros of \texorpdfstring{$\zeta$}{zeta}}; cited by Tome~7, Tome~13, Tome~14, Tome~P1.
\item[5] Proposition~6.2, in \emph{The Zeros of \texorpdfstring{$\zeta$}{zeta}}; cited by Tome~13, Tome~14, Tome~S3, Tome~P1.
\item[6] Theorem~4.19 (Stone--Weierstrass), in \emph{Compactness}; cited by Tome~7, Tome~8, Tome~9, Tome~14.
\item[6] Theorem~8.6 (Fubini), in \emph{Product Measures and Change of Variables}; cited by Tome~7, Tome~8, Tome~10, Tome~P2.
\item[1] Theorem~3.5 (Chinese remainder theorem), in \emph{Congruences}; cited by Tome~14, Tome~C2, Tome~C3.
\item[1] Theorem~4.4 (M\"obius inversion), in \emph{Arithmetic Functions}; cited by Tome~4, Tome~5, Tome~C1.
\item[1] Proposition~4.5 (Gauss), in \emph{Arithmetic Functions}; cited by Tome~4, Tome~C1, Tome~P1.
\item[1] Theorem~5.3 (Gauss), in \emph{Primitive Roots and the Multiplicative Group}; cited by Tome~4, Tome~C2, Tome~S3.
\item[1] Theorem~5.5, in \emph{Primitive Roots and the Multiplicative Group}; cited by Tome~12, Tome~14, Tome~C2.
\item[1] Proposition~8.1, in \emph{Continued Fractions and Pell's Equation}; cited by Tome~10, Tome~P1, Tome~P2.
\item[1] Theorem~8.2, in \emph{Continued Fractions and Pell's Equation}; cited by Tome~5, Tome~10, Tome~P2.
\item[2] Theorem~10.6 (fundamental theorem of Galois theory), in \emph{Galois Theory}; cited by Tome~9, Tome~12, Tome~S3.
\item[2] Proposition~10.9, in \emph{Galois Theory}; cited by Tome~7, Tome~9, Tome~12.
\item[2] Theorem~6.9 (Chinese remainder theorem for rings), in \emph{Rings and Ideals}; cited by Tome~9, Tome~12, Tome~C2.
\item[2] Theorem~7.6 (structure of finitely generated modules over a PID), in \emph{Modules over a Principal Ideal Domain}; cited by Tome~9, Tome~14, Tome~15.
\end{results}

